\documentclass[11pt,a4paper]{article}
\usepackage[T1]{fontenc}
\usepackage{lmodern}
\usepackage[margin=24mm]{geometry}
\usepackage{amsmath,amssymb}
\usepackage{microtype}
\usepackage{fancyhdr,lastpage}
\usepackage[hidelinks]{hyperref}
\hypersetup{pdftitle={Proofs and corrections of OEIS conjectures: overview},pdfauthor={Dev D. Goyal}}
\setlength{\parindent}{0pt}
\setlength{\parskip}{4pt}
\raggedbottom
\pagestyle{fancy}
\fancyhf{}
\fancyhead[L]{\footnotesize\scshape Proofs and corrections of OEIS conjectures}
\fancyhead[R]{\footnotesize\scshape Dev D. Goyal}
\fancyfoot[R]{\footnotesize\thepage\,/\,\pageref*{LastPage}}
\renewcommand{\headrulewidth}{0.3pt}
\fancypagestyle{firstpage}{\fancyhf{}\renewcommand{\headrulewidth}{0pt}\fancyfoot[R]{\footnotesize\thepage\,/\,\pageref*{LastPage}}}
\newcommand{\entry}[5]{\par\medskip\noindent\begin{minipage}{\linewidth}\raggedright
{\bfseries #1.}\enspace #2\par\vspace{2pt}
{\small #3\enspace\emph{(#4.)}\par}
\vspace{1pt}{\footnotesize\href{https://github.com/dvgyl/oeis-a396843-proof/blob/main/preprints/#5/paper.pdf}{Paper}
\enspace\textperiodcentered\enspace\href{https://oeis.org/#1}{OEIS #1}\enspace\textperiodcentered\enspace
\href{https://github.com/dvgyl/oeis-a396843-proof/blob/main/lean/Entries/#1.lean}{Lean}\par}
\end{minipage}\par}
\begin{document}
\thispagestyle{firstpage}
{\LARGE\bfseries Proofs and corrections of OEIS conjectures\par}
\vspace{4pt}
{\large 29 papers in 7 families\par}
\vspace{6pt}
Dev D. Goyal \enspace\textperiodcentered\enspace October 2026
\par\bigskip
This collection settles conjectures recorded in 29 entries of the On-Line Encyclopedia of Integer Sequences (OEIS).
Most of the sequences are coefficients of formal power series defined by functional equations, and most of the
conjectures concern their residues modulo small integers. In most entries the conjectured statements are proved. In
four entries (A380680, A396798, A396843 and A396844) a conjectured statement is false; the papers give
counterexamples and prove the correct statements. Each paper is self-contained, and
every theorem in it is formalized in Lean~4 with Mathlib; the formal proofs use only Lean's standard axioms. All
statements were also checked numerically against the published OEIS data.

The results were obtained in an autonomous AI research effort directed by Dev~D.~Goyal: ChatGPT (OpenAI, GPT-6 with
Pro reasoning) found the proofs and wrote the Lean formalization from Dev~D.~Goyal's prompts, and Claude Opus~5.5,
used with Claude Code, checked the work independently and wrote the papers in their present form. Dev~D.~Goyal has
verified the results. The papers have not been peer reviewed.
\section*{1. The family $A_m(xA_m-mxA_m^2)=x^2$}
Parity, residues modulo $m^2$ and binary-digit divisibility for the unique normalized solution of $A_m(xA_m-mxA_m^2)=x^2$, with corrected residue classifications for the two entries.
\entry{A396843}{Odd coefficients, remainders modulo $9$, and binary digits for the series $A(xA-3xA^2)=x^2$}{Every coefficient is odd; the exact residues modulo 9 (correcting the entry's second conjecture); $3^{s_2(n-1)}$ divides $a(n)$, where $s_2$ is the binary digit sum.}{proved; one statement corrected}{A396843-Odd-coefficients-remainders-modulo-9-and-binary-digits}
\entry{A396844}{Remainders modulo $16$ for the series $A(xA-4xA^2)=x^2$}{The exact residues modulo 16; the conjecture fails first at $n=17$.}{corrected}{A396844-Remainders-modulo-16}
\section*{2. Consecutive iterates $F=x+F^{\circ r}\,F^{\circ(r+1)}$}
For every integer $k$, the $k$-th iterate of $F_r$ modulo $2r$ is $x/(1-kx)$ when $r$ is odd, with an explicit correction term when $r$ is even.
\entry{A396099}{The iterates of a self-referential power series modulo $4$}{All four conjectured congruences modulo 4 for the series and its iterates.}{proved}{A396099-Iterates-of-a-self-referential-series-modulo-4}
\entry{A396798}{The iterates of a self-referential power series modulo $8$}{Seven of the eight conjectured congruences modulo 8; the remaining one holds from $n=2$, not from $n=1$ as stated.}{proved; one statement corrected}{A396798-Iterates-of-a-self-referential-series-modulo-8}
\entry{A396807}{The iterates of a self-referential power series modulo $10$}{For every integer $k$, the coefficients of the $k$-th iterate satisfy $[x^n]\,A^{\circ k}\equiv k^{n-1}\pmod{10}$.}{proved}{A396807-Iterates-of-a-self-referential-series-modulo-10}
\section*{3. Scaled inverses $A(x-x\,A(qx)/q)=x$}
Every coefficient is $1$ modulo $q+1$; for even $q$, the odd coefficients sit exactly at the powers of $2$.
\entry{A393856}{Congruences for the series $A$ with $A\bigl(x-x\,A(4x)/4\bigr)=x$}{$a(n)\equiv1\pmod 5$, and $a(n)$ is odd exactly when $n$ is a power of 2.}{proved}{A393856-Congruences-for-a-scaled-inverse-series-q-4}
\entry{A393857}{Congruences for the series $A$ with $A\bigl(x-x\,A(5x)/5\bigr)=x$}{$a(n)\equiv1\pmod 6$.}{proved}{A393857-Congruences-for-a-scaled-inverse-series-q-5}
\entry{A393858}{Congruences for the series $A$ with $A\bigl(x-x\,A(6x)/6\bigr)=x$}{$a(n)\equiv1\pmod 7$, and $a(n)$ is odd exactly when $n$ is a power of 2.}{proved}{A393858-Congruences-for-a-scaled-inverse-series-q-6}
\entry{A393859}{Congruences for the series $A$ with $A\bigl(x-x\,A(7x)/7\bigr)=x$}{$a(n)\equiv1\pmod 8$.}{proved}{A393859-Congruences-for-a-scaled-inverse-series-q-7}
\section*{4. Theta series and theta-type equations}
Series built from the Jacobi theta series $\theta_3$ or defined by theta-type sums, whose coefficients follow simple patterns modulo 4 or 16.
\entry{A378580}{Remainders modulo $4$ for the series with $A\bigl(x/A(x)\bigr)=\theta_3(x)$}{$A(x/A(x))=\theta_3$: for $n>0$, $a(n)\equiv2\pmod4$ if $n$ is a square, and $4\mid a(n)$ otherwise.}{proved}{A378580-Theta-quotient-substitution-modulo-4}
\entry{A378581}{Remainders modulo $4$ for the series with $A\bigl(x\,A(x)\bigr)=\theta_3(x)$}{$A(x\,A(x))=\theta_3$: the same square pattern.}{proved}{A378581-Theta-product-substitution-modulo-4}
\entry{A380060}{Square indices and remainders modulo 4 in a theta-type series}{For $n>0$, $a(n)\equiv2\pmod4$ if $n$ is a square, and $4\mid a(n)$ otherwise.}{proved}{A380060-Square-indices-in-a-theta-type-series}
\entry{A380680}{Twice the squares, not the squares: remainders modulo 4 in a theta-type series}{The square conjecture is false ($a(1)=4$, $a(2)=6$); in fact $a(n)\equiv2\pmod4$ exactly when $n=2k^2$.}{conjecture disproved; corrected statement proved}{A380680-Twice-the-squares-not-the-squares}
\entry{A355872}{Residues modulo $16$ for a series defined by a bilateral theta-type equation}{$a(n)\equiv2(-1)^{n-1}\pmod{16}$ for all $n\ge1$, which implies the three conjectured congruences modulo 4 and 8.}{proved}{A355872-Residues-modulo-16-for-a-bilateral-theta-type-equation}
\section*{5. Catalan-type equations}
Equations whose solutions reduce, modulo 4 or 8, to expressions in the Catalan series, whose coefficients are odd exactly at powers of 2.
\entry{A393172}{The Catalan numbers inside the equation $A(x-x^2)=x+A(x)^2$}{For $n>1$, $a(n)\equiv2\pmod4$ if $n=2^k$, and $4\mid a(n)$ otherwise.}{proved}{A393172-The-Catalan-numbers-inside-a-functional-equation}
\entry{A393171}{Remainders modulo $8$ of a series built from the Catalan numbers}{For $n>2$, $a(n)\equiv4\pmod8$ if $n=2^k$ or $2^k+1$, and $8\mid a(n)$ otherwise.}{proved}{A393171-Remainders-modulo-8-from-the-Catalan-numbers}
\entry{A392210}{The odd-indexed terms of a nested square are even}{$a(2n-1)$ is even for $n>1$ (an elementary observation recorded in the entry).}{proved}{A392210-Odd-indexed-terms-of-a-nested-square}
\entry{A389536}{Remainders modulo $4$ for the equation $A(x^2+2x^3)=x\,A(x)$}{For $n>1$, $a(n)\equiv2\pmod4$ if $n=2^k+1$, and $4\mid a(n)$ otherwise.}{proved}{A389536-Remainders-modulo-4-for-a-Mahler-type-equation}
\entry{A397588}{Odd coefficients at powers of two for $A=x+(xA^2)'$}{$a(n)$ is odd exactly when $n$ is a power of 2, and $3\mid a(n)$ for $n>1$.}{proved}{A397588-Odd-coefficients-at-powers-of-two}
\entry{A195194}{Parity of an infinitely nested series reversion}{The infinitely nested reversion is well defined, and $a(n)$ is odd exactly when $n$ is a power of 2.}{proved}{A195194-Parity-of-an-infinitely-nested-series-reversion}
\section*{6. Reversions with an exact halving step}
Equations whose coefficient recursion appears to divide by 2; the division is shown to be exact before reducing.
\entry{A389537}{Congruences modulo $4$ for the series with $A\bigl(A(x)-x\bigr)=A(x)^2$}{For $n>2$, $a(n)\equiv2\pmod4$ if $n=3\cdot2^k$, and $4\mid a(n)$ otherwise.}{proved}{A389537-Congruences-modulo-4-for-a-sparse-reversion}
\entry{A389542}{Congruences modulo $4$ for the series with $A\bigl(A(x)^2-x^2\bigr)=4A(x)^3$}{For $n>1$, $a(n)\equiv2\pmod4$ if $n=2^k+1$, and $4\mid a(n)$ otherwise.}{proved}{A389542-Congruences-modulo-4-for-a-square-reversion}
\entry{A389474}{A power series whose square is $x^2+2x^3$ modulo $9$}{$A^2\equiv x^2+2x^3\pmod 9$, and $A^{3k}\equiv A(x^3)^k\pmod 3$.}{proved}{A389474-A-series-whose-square-is-a-polynomial-modulo-9}
\entry{A389475}{Two congruences modulo $3$ for the series with $A(3x^2)=2x^2+A^2-2A^3$}{$A^3\equiv A(x^3)$ and $A^2\equiv x^2+2A(x^3)\pmod 3$.}{proved}{A389475-Two-congruences-modulo-3}
\section*{7. Other results}
Further entries, each with its own method.
\entry{A396093}{When are the coefficients of $B(B(B(x)))$ odd?}{$a(n)$ is odd exactly when $n\equiv1,3,7,9\pmod{10}$.}{proved}{A396093-Odd-coefficients-of-the-third-iterate}
\entry{A396102}{Every coefficient is $1$ modulo $3$: the series with $A(A(A(x)))=(1+x)\,A(A(x))$}{$a(n)\equiv1\pmod 3$, together with a theorem for the whole family $F^{\circ r}=(1+x)F^{\circ(r-1)}$.}{proved}{A396102-Every-coefficient-is-1-modulo-3}
\entry{A380056}{The coefficients of $(e^x-1)/\cos 2x$ modulo 4 and 5}{$a(n)\equiv1\pmod4$, and $a(n)\equiv2\cdot2^n+3\cdot4^n\pmod5$, so $5\mid a(n)$ exactly when $4\mid n$.}{proved}{A380056-Exponential-over-cosine-modulo-4-and-5}
\entry{A397346}{Residues modulo $8$ for a series defined by a self-referential exponential}{From $n=2$ on, $a(n)\bmod 8$ repeats $4,2,0,2$; more precisely, $a(n)/n\equiv6\pmod8$ if $n\equiv3\pmod4$ and $a(n)/n\equiv2\pmod8$ otherwise.}{proved}{A397346-Residues-modulo-8-for-a-self-referential-exponential}
\entry{A394576}{A greedy sequence that avoids sums of three consecutive terms}{$\lfloor(4n-3)/3\rfloor\le a(n)\le\lfloor(4n-1)/3\rfloor$ for all $n\ge1$; in particular $a(3k)=4k-1$.}{proved}{A394576-A-greedy-sequence-avoiding-sums-of-three-consecutive-terms}
\end{document}
